% Part: sets-functions-relations
% Chapter: size-of-sets-complete

\documentclass[../../../include/open-logic-chapter]{subfiles}

\begin{document}

\olchapter{sfr}{siz}{د سټونو اندازه}

\begin{editorial}
دا باب شمېرنې، د شمېر وړ والے او د شمېر نۀ وړ والے څېړي۔
ځينې برخې دوه بڼې لري: يوه زياته ابتدايي بڼه، چې شمېرنې د
لړونو يا له $\PosInt$ څخه د هغو تابعو په توګه اخلي چې پر هدف
سټ بشپړې وي؛ او يوه زياته تجريدي بڼه، چې شمېرنې له $\Nat$ سره
د دوه اړخيزو يو پر يو تابعو په توګه تعريفوي۔
\end{editorial}

\olimport{introduction}

\olimport{enumerability}

\olimport{zig-zag}

\olimport{pairing}

\olimport{pairing-alt}

\olimport{non-enumerability}

\olimport{reduction}

\olimport{equinumerous-sets}

\olimport{comparing-size}

\olimport{schroder-bernstein}

\begin{editorial}
لاندې \olref[sfr][siz][enm-alt]{sec}،
\olref[sfr][siz][nen-alt]{sec} او \olref[sfr][siz][red-alt]{sec}
د \olref[sfr][siz][enm]{sec}، \olref[sfr][siz][nen]{sec} او
\olref[sfr][siz][red]{sec} بديلې بڼې دي چې ټيم بټن د خپل
«د سټونو خلاصه نظريه» متن کښې د کارولو دپاره برابرې کړې دي۔
دوې لږې زياتې پرمختللې دي او د شمېر وړ والي يو بېل تعريف
کاروي چې د سټونو د نظريې په سياق کښې زيات مناسب دے؛ يعنې
له $\Nat$ يا د هغې له پيلنۍ برخې سره دوه اړخيزه يو پر يو تابع
لرل، د دې پر ځاے چې په لړ کښې د راوړلو وړ وي يا له $\PosInt$
څخه د يوې پر هدف سټ بشپړې تابعې د قيمتونو سټ وي۔
\end{editorial}

\olimport{enumerability-alt}
\olimport{non-enumerability-alt}
\olimport{reduction-alt}

\OLEndChapterHook

\end{document}
